%=============================================================
%  ARTÍCULO DE REVISIÓN SISTEMÁTICA DE LITERATURA
%  Aplicación de ML en Sistemas de Recomendación Multiagente
%  para la Optimización de Ventas en E-Commerce
%  Protocolo: Kitchenham & Brereton (2007)
%  Escuela Profesional de Ingeniería de Sistemas — UPeU Juliaca
%=============================================================

\documentclass[12pt, a4paper]{article}

\PassOptionsToPackage{table, xcdraw}{xcolor}
\usepackage[utf8]{inputenc}
\usepackage[T1]{fontenc}
\usepackage[spanish, es-nodecimaldot]{babel}
\usepackage{csquotes}
\usepackage{lmodern}
\usepackage{microtype}
\usepackage{setspace}
\onehalfspacing
\usepackage[top=3cm, bottom=2.5cm, left=3cm, right=2.5cm]{geometry}
\usepackage{amsmath, amssymb}
\usepackage{graphicx}
\usepackage{float}
\usepackage{caption}
\usepackage{subcaption}
\usepackage{tikz}
\usetikzlibrary{shapes.geometric, arrows.meta, positioning,
                fit, backgrounds, calc}
\usepackage{booktabs}
\usepackage{tabularx}
\usepackage{array}
\usepackage{multirow}
\usepackage{longtable}
\renewcommand{\arraystretch}{1.35}
\usepackage{enumitem}
\setlist[itemize]{leftmargin=1.5em, itemsep=2pt}
\setlist[enumerate]{leftmargin=1.5em, itemsep=2pt}
\usepackage{xcolor}
\definecolor{azulPrincipal}{RGB}{0, 70, 127}
\definecolor{azulMedio}{RGB}{0, 102, 179}
\definecolor{azulClaro}{RGB}{204, 229, 255}
\definecolor{grisClaro}{RGB}{245, 245, 245}
\definecolor{verdeOsc}{RGB}{0, 100, 60}
\definecolor{naranjaAlerta}{RGB}{255, 140, 0}
\definecolor{rojoAlerta}{RGB}{180, 0, 0}

\usepackage{tcolorbox}
\tcbuselibrary{skins, breakable, theorems}

\newtcolorbox{preguntabox}[1][]{
  enhanced, colback=azulClaro!50, colframe=azulPrincipal,
  fonttitle=\bfseries\color{white}, coltitle=white,
  colbacktitle=azulPrincipal,
  title={Preguntas de Investigaci\'on}, #1}

\newtcolorbox{notamet}{
  enhanced, breakable, colback=naranjaAlerta!15,
  colframe=naranjaAlerta, leftrule=4pt, left=8pt,
  fontupper=\small\itshape}

\newtcolorbox{hallazgo}[1][]{
  enhanced, colback=verdeOsc!8, colframe=verdeOsc,
  fonttitle=\bfseries\color{white}, coltitle=white,
  colbacktitle=verdeOsc, title={Hallazgo}, #1}

\usepackage{fancyhdr}
\setlength{\headheight}{15pt}
\pagestyle{fancy}
\fancyhf{}
\fancyhead[L]{\small\textcolor{azulPrincipal}{\textit{RSL --- ML en Sistemas de Recomendaci\'on Multiagente para E-Commerce}}}
\fancyhead[R]{\small\textcolor{azulPrincipal}{\textit{Art\'iculo de Revisi\'on}}}
\fancyfoot[C]{\small\thepage}
\renewcommand{\headrulewidth}{0.4pt}

\usepackage{titlesec}
\titleformat{\section}
  {\large\bfseries\color{azulPrincipal}}{\thesection.}{0.6em}{}[\titlerule]
\titleformat{\subsection}
  {\normalsize\bfseries\color{azulMedio}}{\thesubsection.}{0.5em}{}
\titleformat{\subsubsection}
  {\normalsize\itshape\color{azulPrincipal}}{\thesubsubsection.}{0.4em}{}

\usepackage{hyperref}
\hypersetup{
  colorlinks=true, linkcolor=azulPrincipal,
  citecolor=azulMedio, urlcolor=azulMedio,
  pdftitle={RSL ML Sistemas Recomendacion Multiagente E-Commerce},
}

\usepackage[backend=biber, style=ieee, sorting=nyt, maxbibnames=6]{biblatex}
\addbibresource{referencias.bib}

\newcommand{\RSL}{Revisi\'on Sistem\'atica de Literatura}

%=============================================================
\begin{document}
%=============================================================

%-------------------------------------------------------------
% PORTADA
%-------------------------------------------------------------
\begin{center}
  {\LARGE\bfseries\color{azulPrincipal}
    Art\'iculo de revisi\'on sistem\'atica de literatura sobre
    la aplicaci\'on de t\'ecnicas de Machine Learning en sistemas
    de recomendaci\'on basados en arquitecturas multiagente para
    la optimizaci\'on de ventas en plataformas de comercio
    electr\'onico}\\[1em]

  {\large
    Nombre Apellido\textsuperscript{1},\quad
    Nombre Apellido\textsuperscript{2},\quad
    Nombre Apellido\textsuperscript{3}}\\[0.5em]

  {\small
    \textsuperscript{1}Escuela Profesional de Ingenier\'ia de Sistemas,
    Universidad Peruana Uni\'on, Juliaca, Per\'u.
    \href{mailto:autor1@upeu.edu.pe}{autor1@upeu.edu.pe}\\
    \textsuperscript{2}Escuela Profesional de Ingenier\'ia de Sistemas,
    Universidad Peruana Uni\'on, Juliaca, Per\'u.
    \href{mailto:autor2@upeu.edu.pe}{autor2@upeu.edu.pe}\\
    \textsuperscript{3}Escuela Profesional de Ingenier\'ia de Sistemas,
    Universidad Peruana Uni\'on, Juliaca, Per\'u.
    \href{mailto:autor3@upeu.edu.pe}{autor3@upeu.edu.pe}}\\[0.8em]

  {\small\textbf{Recibido:} DD/MM/2025\quad
         \textbf{Aceptado:} DD/MM/2025\quad
         \textbf{Publicado:} DD/MM/2025}
\end{center}

\vspace{0.5em}\noindent\rule{\linewidth}{1pt}

%-------------------------------------------------------------
% RESUMEN
%-------------------------------------------------------------
\section*{Resumen}
\noindent
\textbf{Antecedentes:} El comercio electr\'onico ha experimentado un
crecimiento exponencial impulsando la demanda de sistemas de recomendaci\'on
precisos y personalizados. La integraci\'on de t\'ecnicas de Machine Learning
(ML) con arquitecturas multiagente emerge como enfoque prometedor para
mejorar la experiencia del usuario y optimizar ventas en plataformas digitales.
\textbf{Objetivo:} Identificar, analizar y sintetizar la evidencia cient\'ifica
sobre la aplicaci\'on de t\'ecnicas de ML en sistemas de recomendaci\'on basados
en arquitecturas multiagente orientados a la optimizaci\'on de ventas en
plataformas de e-commerce, publicada entre 2020 y 2026.
\textbf{M\'etodo:} Se realiz\'o una \RSL{} siguiendo el protocolo de
Kitchenham \& Brereton (2007), consultando Scopus, IEEE Xplore, ACM Digital
Library, SpringerLink y ScienceDirect. Se aplicaron nueve criterios de
inclusi\'on y ocho de exclusi\'on con apoyo de la herramienta Parsif.al.
\textbf{Resultados:} Se seleccionaron 31 estudios primarios. Los hallazgos
evidencian que el Deep Learning, el filtrado colaborativo neuronal y el
aprendizaje por refuerzo son las t\'ecnicas de ML m\'as utilizadas. Las
arquitecturas multiagente mejoran la escalabilidad y la precisi\'on de las
recomendaciones de forma significativa.
\textbf{Conclusiones:} La combinaci\'on de ML con arquitecturas multiagente
representa el estado del arte en sistemas de recomendaci\'on para e-commerce,
aunque persisten brechas en interoperabilidad, privacidad y evaluaci\'on en
entornos reales.

\vspace{0.4em}
\noindent\textbf{Palabras clave:} machine learning, sistemas de recomendaci\'on,
arquitecturas multiagente, comercio electr\'onico, deep learning, optimizaci\'on
de ventas, aprendizaje por refuerzo.

\vspace{0.5em}\noindent\rule{\linewidth}{0.4pt}

%-------------------------------------------------------------
% ABSTRACT
%-------------------------------------------------------------
\section*{Abstract}
\noindent
\textbf{Background:} E-commerce has experienced exponential growth driving
the demand for accurate and personalized recommendation systems. The
integration of Machine Learning (ML) techniques with multi-agent architectures
emerges as a promising approach to improve user experience and optimize sales.
\textbf{Objective:} To identify, analyze and synthesize scientific evidence on
the application of ML techniques in recommendation systems based on multi-agent
architectures aimed at optimizing sales in e-commerce platforms, published
between 2020 and 2026.
\textbf{Method:} A Systematic Literature Review (SLR) was conducted following
the Kitchenham \& Brereton (2007) protocol, consulting Scopus, IEEE Xplore,
ACM Digital Library, SpringerLink and ScienceDirect. Nine inclusion and eight
exclusion criteria were applied using the Parsif.al tool.
\textbf{Results:} 31 primary studies were selected. Findings show that Deep
Learning, neural collaborative filtering and reinforcement learning are the most
widely used ML techniques. Multi-agent architectures significantly improve
scalability and accuracy of recommendations.
\textbf{Conclusions:} The combination of ML with multi-agent systems represents
the state-of-the-art in e-commerce recommendation systems, although gaps remain
regarding interoperability, privacy and evaluation in real environments.

\vspace{0.4em}
\noindent\textbf{Keywords:} machine learning, recommender systems, multi-agent
architectures, e-commerce, deep learning, sales optimization, reinforcement
learning.

\noindent\rule{\linewidth}{1pt}\vspace{1em}

%=============================================================
\section{Introducci\'on}
%=============================================================

El comercio electr\'onico se ha consolidado como uno de los sectores de mayor
dinamismo en la econom\'ia digital global. Plataformas como Amazon, Alibaba
y Taobao generan millones de transacciones diarias en las que los sistemas de
recomendaci\'on desempe\~nan un papel determinante: se estima que alrededor
del 35\,\% de los ingresos de Amazon son atribuibles directamente a su motor
de recomendaciones \cite{Girish2025}. En este contexto, la capacidad de
personalizar las sugerencias de productos en tiempo real constituye una ventaja
competitiva cr\'itica para cualquier plataforma digital \cite{Halachev2024}.

Los sistemas de recomendaci\'on tradicionales, basados en filtrado colaborativo
y filtrado basado en contenido, enfrentan limitaciones ante la escala y la
dinamicidad de los entornos actuales: problemas de arranque en fr\'io
(\textit{cold start}), escasez de datos (\textit{data sparsity}) y dificultad
para capturar preferencias no lineales del usuario \cite{Loukili2023,
Alam2025}. La irrupci\'on del Machine Learning (ML) y, en particular, del
aprendizaje profundo (\textit{deep learning}), ha permitido superar varias de
estas barreras mediante la extracci\'on autom\'atica de representaciones
latentes a partir de datos de interacci\'on usuario-\'item \cite{Nagraj2024,
Lampa2024}.

Paralelamente, las arquitecturas multiagente han emergido como un paradigma
arquitect\'onico capaz de distribuir, coordinar y escalar los procesos de
recomendaci\'on mediante agentes inteligentes aut\'onomos que colaboran para
generar sugerencias m\'as precisas y contextuales \cite{Alhejaili2021,
Zhang2025DARLR}. La combinaci\'on de ambos enfoques ---ML y sistemas
multiagente--- est\'a siendo explorada activamente en la literatura acad\'emica,
aunque no existe a\'un una s\'intesis sistem\'atica que integre sus hallazgos,
m\'etodos y brechas de investigaci\'on.

La presente \RSL{} responde a esa necesidad, siguiendo el protocolo de
Kitchenham \& Brereton \cite{Kitchenham2007}, est\'andar reconocido en
ingenier\'ia de software y ciencias de la computaci\'on para la s\'intesis de
evidencia emp\'irica. Las contribuciones de este art\'iculo son: (i)~una
s\'intesis exhaustiva de las t\'ecnicas de ML empleadas en sistemas de
recomendaci\'on multiagente para e-commerce; (ii)~un an\'alisis comparativo de
las m\'etricas de evaluaci\'on; (iii)~la identificaci\'on de brechas de
investigaci\'on abiertas; y (iv)~un mapa de evidencia que relaciona los estudios
primarios con las preguntas de investigaci\'on. El art\'iculo se organiza as\'i:
\S\,2 presenta el marco conceptual; \S\,3 detalla el proceso metodol\'ogico;
\S\,4 reporta los resultados; \S\,5 discute los hallazgos; \S\,6 formula las
conclusiones.

\begin{preguntabox}
\begin{enumerate}[label=\textbf{RQ\arabic*.}, leftmargin=2em]
  \item \textquestiondown Qu\'e t\'ecnicas de Machine Learning aplicadas en
        sistemas de recomendaci\'on basados en arquitecturas multiagente han
        demostrado mayor efectividad en la optimizaci\'on de ventas en
        plataformas de comercio electr\'onico?
  \item \textquestiondown C\'omo contribuyen las arquitecturas multiagente a
        mejorar el rendimiento de los sistemas de recomendaci\'on basados en
        Machine Learning?
  \item \textquestiondown Qu\'e m\'etricas se emplean para evaluar la
        efectividad de los sistemas de recomendaci\'on en t\'erminos de
        precisi\'on, conversi\'on y ventas?
  \item \textquestiondown C\'omo se comparan los sistemas de recomendaci\'on
        tradicionales con aquellos basados en ML y arquitecturas multiagente en
        el contexto del comercio electr\'onico?
  \item \textquestiondown Cu\'al es el impacto de estos sistemas en el
        comportamiento del usuario y en la optimizaci\'on de ventas?
  \item \textquestiondown Cu\'ales son los principales desaf\'ios, limitaciones
        y brechas de investigaci\'on en la implementaci\'on de estos sistemas en
        entornos reales de e-commerce?
\end{enumerate}
\end{preguntabox}

%=============================================================
\section{Marco Conceptual}
%=============================================================

\subsection{Machine Learning en sistemas de recomendaci\'on}

El Machine Learning (ML) es la subdisciplina de la inteligencia artificial que
dota a los sistemas de la capacidad de aprender patrones a partir de datos sin
ser expl\'icitamente programados \cite{Halachev2024}. En el contexto de los
sistemas de recomendaci\'on para e-commerce, las t\'ecnicas de ML m\'as
relevantes son:

\textbf{Aprendizaje profundo (\textit{Deep Learning}):} Emplea redes neuronales
multicapa para extraer representaciones jer\'arquicas. Sus variantes m\'as
usadas incluyen redes convolucionales (CNN), redes recurrentes (RNN/LSTM),
mecanismos de atenci\'on (\textit{Transformer}) y modelos de memoria din\'amica
\cite{Gupta2024, Ding2025}. Los sistemas basados en Deep Learning superan
sistem\'aticamente a los m\'etodos tradicionales en precisi\'on y recuperaci\'on
\cite{Alam2025, Lampa2024}.

\textbf{Filtrado colaborativo neuronal y factorizaci\'on matricial:} Extienden
el filtrado colaborativo cl\'asico mediante redes neuronales, modelando
interacciones no lineales usuario-\'item \cite{LiYuetong2025, Girish2025}.
La factorizaci\'on matricial con variantes profundas maneja eficientemente
datos dispersos a gran escala \cite{Lin2025}.

\textbf{Aprendizaje por refuerzo (\textit{Reinforcement Learning}):} Modela
la recomendaci\'on como un proceso de decisi\'on de Markov (MDP), optimizando
pol\'iticas de recomendaci\'on a largo plazo \cite{Zhang2025DARLR,
Zhang2023RL}. El aprendizaje por refuerzo offline y el marco de doble agente
DARLR son propuestas de frontera \cite{Zhang2025DARLR}.

\textbf{Modelos h\'ibridos y ensemble:} La combinaci\'on de filtrado
colaborativo, filtrado basado en contenido y Deep Learning maximiza cobertura
y precisi\'on \cite{Gupta2024, Lu2025, Girish2025}. Los clasificadores ensemble
como XGBoost y Random Forest demuestran alta robustez ante datos desbalanceados
\cite{Dinata2025}.

\subsection{Arquitecturas multiagente en e-commerce}

Un sistema multiagente (SMA) est\'a compuesto por m\'ultiples agentes
inteligentes aut\'onomos que interact\'uan en un entorno compartido para
alcanzar objetivos individuales o colectivos \cite{Alhejaili2021}. En
e-commerce, los SMA permiten distribuir el proceso de recomendaci\'on entre
agentes especializados: agentes de perfil de usuario, agentes de an\'alisis de
\'items, agentes de negociaci\'on y agentes de presentaci\'on \cite{Alhejaili2021,
Bai2025}.

Las ventajas estructurales de los SMA sobre los sistemas monol\'iticos incluyen:
(i)~escalabilidad horizontal mediante paralelizaci\'on; (ii)~modularidad que
facilita actualizaciones incrementales; (iii)~robustez ante fallos parciales;
y (iv)~adaptaci\'on din\'amica al contexto del usuario \cite{Ho2025GraphMind,
Bai2025}. En aplicaciones recientes, los SMA han incorporado Modelos de
Lenguaje de Gran Escala (LLM) y redes de atenci\'on de grafos (GAE) para
gestionar el contexto conversacional de forma eficiente \cite{Ho2025GraphMind}.

\subsection{Sinergia entre ML, multiagente y optimizaci\'on de ventas}

La sinergia entre ML y arquitecturas multiagente en e-commerce opera en tres
niveles: (i)~a nivel de \textit{datos}, los agentes recopilan y procesan datos
multimodales de comportamiento del usuario que alimentan modelos de ML
especializados; (ii)~a nivel de \textit{inferencia}, diferentes agentes ejecutan
distintos modelos de recomendaci\'on cuyas salidas son coordinadas mediante
protocolos de negociaci\'on o enrutamiento inteligente \cite{Alhejaili2021,
Bai2025}; y (iii)~a nivel de \textit{negocio}, la optimizaci\'on conjunta de
m\'etricas de precisi\'on y conversi\'on requiere considerar tanto la relevancia
de las recomendaciones como su impacto en el comportamiento de compra
\cite{Girish2025, Ding2025, LiYuetong2025}.

%=============================================================
\section{Proceso Metodol\'ogico}
%=============================================================

Esta revisi\'on sigue el protocolo de Kitchenham \& Brereton
\cite{Kitchenham2007} para revisiones sistem\'aticas en ingenier\'ia de
software, adaptado al \'area de ML e inteligencia artificial aplicada al
comercio electr\'onico, y ejecutado con la plataforma Parsif.al
(\url{https://parsif.al}).

\subsection{Protocolo de Revisi\'on (PICOC)}

\begin{table}[H]
\centering
\caption{Marco PICOC de la revisi\'on sistem\'atica.}
\label{tab:picoc}
\begin{tabularx}{\linewidth}{>{\bfseries}lX}
\toprule
Elemento & Descripci\'on \\
\midrule
Poblaci\'on &
  Plataformas de comercio electr\'onico y usuarios que interact\'uan con
  sistemas de recomendaci\'on basados en Machine Learning y arquitecturas
  multiagente. \\
Intervenci\'on &
  Aplicaci\'on de t\'ecnicas de Machine Learning en sistemas de
  recomendaci\'on basados en arquitecturas multiagente. \\
Comparaci\'on &
  Sistemas de recomendaci\'on tradicionales sin uso de Machine Learning
  ni arquitecturas multiagente. \\
Resultado &
  Mejora en la precisi\'on de las recomendaciones, optimizaci\'on de
  ventas y mejora en la experiencia del usuario en plataformas de
  comercio electr\'onico. \\
Contexto &
  Plataformas de comercio electr\'onico (e-commerce). \\
\bottomrule
\end{tabularx}
\end{table}

\subsection{Bases de Datos y Fuentes}

Se consultaron las siguientes fuentes electr\'onicas de literatura acad\'emica
indexada:
\begin{itemize}
  \item Scopus (\url{https://www.scopus.com})
  \item IEEE Xplore (\url{https://ieeexplore.ieee.org})
  \item ACM Digital Library (\url{https://dl.acm.org})
  \item SpringerLink (\url{https://link.springer.com})
  \item ScienceDirect / Elsevier (\url{https://www.sciencedirect.com})
  \item Parsif.al --- herramienta de apoyo RSL (\url{https://parsif.al})
\end{itemize}

\subsection{Cadena de B\'usqueda}

La cadena fue construida a partir de los t\'erminos PICOC y sus sin\'onimos,
agrupados en tres bloques tem\'aticos conectados con operadores booleanos
\texttt{AND}. Se valid\'o en Parsif.al y se adapt\'o a la sintaxis de cada
base de datos.

\begin{notamet}
  La cadena fue registrada en el protocolo de Parsif.al antes de ejecutar
  la b\'usqueda, garantizando trazabilidad y reproducibilidad.
\end{notamet}

\vspace{0.5em}
\noindent\textbf{Cadena de b\'usqueda principal:}

\begin{center}
\fbox{\parbox{0.95\linewidth}{\small\ttfamily\raggedright
(``E-commerce'' OR ``digital commerce'' OR ``digital platforms'' OR
``electronic commerce'' OR ``online retail'' OR ``online shopping'' OR
``online stores'' OR ``online systems'' OR ``web commerce'' OR
``Recommender Systems'' OR ``collaborative filtering'' OR
``content-based filtering'' OR ``hybrid recommender systems'' OR
``personalization systems'' OR ``recommendation engine'' OR
``recommendation systems'' OR ``recommender engine'')\\
AND\\
(``Algorithms and Models'' OR ``deep recommender models'' OR
``matrix factorization'' OR ``neural collaborative filtering'' OR
``optimization algorithms'' OR ``ranking algorithms'' OR
``reinforcement learning models'' OR ``Machine Learning'' OR
``artificial intelligence'' OR ``deep learning'' OR ``ensemble learning''
OR ``neural networks'' OR ``predictive modeling'' OR
``reinforcement learning'' OR ``supervised learning'' OR
``unsupervised learning'' OR ``Multi-agent Systems'' OR ``Multi-agent''
OR ``agent-based modeling'' OR ``agent-based systems'' OR
``agent-oriented systems'' OR ``autonomous agents'' OR
``distributed agents'' OR ``intelligent agents'' OR ``multiagent systems''
OR ``Software Architecture'' OR ``centralized architecture'' OR
``distributed architecture'' OR ``microservices architecture'' OR
``scalable systems'' OR ``service-oriented architecture'' OR
``system architecture'' OR ``system design'')\\
AND\\
(``Metrics'' OR ``F1-score'' OR ``MAE'' OR ``NDCG'' OR ``RMSE'' OR
``accuracy'' OR ``hit rate'' OR ``performance evaluation'' OR
``precision'' OR ``ranking metrics'' OR ``recall'' OR
``Sales Optimization'' OR ``conversion rate'' OR ``purchase intention''
OR ``revenue optimization'' OR ``sales performance'')
}}
\end{center}

\subsection{Criterios de Inclusi\'on y Exclusi\'on}

\begin{table}[H]
\centering
\caption{Criterios de selecci\'on de estudios primarios.}
\label{tab:criterios}
\begin{tabularx}{\linewidth}{lX}
\toprule
\textbf{C\'odigo} & \textbf{Criterio} \\
\midrule
\multicolumn{2}{l}{\cellcolor{azulClaro!40}\textbf{Criterios de Inclusi\'on}} \\
CI-1 & Estudios en bases indexadas (Scopus, IEEE Xplore, ACM,
       ScienceDirect, Springer). \\
CI-2 & Publicados entre 2020 y 2026. \\
CI-3 & Enfocados en sistemas de recomendaci\'on para e-commerce. \\
CI-4 & Que apliquen Machine Learning / Deep Learning y/o arquitecturas
       multiagente. \\
CI-5 & Que incluyan evaluaci\'on cuantitativa con m\'etricas (accuracy,
       recall, F1, conversi\'on, ventas). \\
CI-6 & Que permitan comparaci\'on de enfoques o incluyan baseline. \\
CI-7 & Publicaciones revisadas por pares. \\
CI-8 & Revistas clasificadas en Q1 o Q2 (Scimago Journal Rank) o
       conferencias reconocidas (CORE A, A*). \\
CI-9 & Idioma ingl\'es. \\
\midrule
\multicolumn{2}{l}{\cellcolor{rojoAlerta!10}\textbf{Criterios de Exclusi\'on}} \\
CE-1 & Duplicados. \\
CE-2 & Fuera del rango temporal (antes de 2020 o despu\'es de 2026). \\
CE-3 & Dominios no relacionados con e-commerce (medicina, educaci\'on,
       agricultura, etc.). \\
CE-4 & Sin validaci\'on emp\'irica o sin m\'etricas reportadas. \\
CE-5 & Versiones redundantes (se conserva la m\'as completa/reciente). \\
CE-6 & Sin acceso a texto completo. \\
CE-7 & Estudios puramente te\'oricos. \\
CE-8 & Revistas fuera de Q1--Q2 o no indexadas (salvo conferencias
       relevantes). \\
\bottomrule
\end{tabularx}
\end{table}

\subsection{Proceso de Selecci\'on}

El proceso de selecci\'on se realiz\'o en cuatro pasos siguiendo las fases
del protocolo Kitchenham \& Brereton \cite{Kitchenham2007}:

\begin{enumerate}
  \item \textbf{Paso 1 --- B\'usqueda inicial (CI-1, CI-2, CI-9):}
        Ejecuci\'on de la cadena en todas las bases de datos; filtro
        autom\'atico por idioma ingl\'es y rango temporal 2020--2026.
        Registro en Parsif.al.
  \item \textbf{Paso 2 --- Cribado por t\'itulo y eliminaci\'on de duplicados
        (CE-1, CE-2, CE-6):}
        Revisi\'on de t\'itulos; eliminaci\'on de duplicados inter-bases y
        art\'iculos sin acceso al texto completo.
  \item \textbf{Paso 3 --- Selecci\'on por resumen (CI-3, CI-4, CE-3, CE-7):}
        Lectura de abstracts; exclusi\'on de estudios fuera del dominio de
        e-commerce o que no emplean ML/multiagente; exclusi\'on de estudios
        puramente te\'oricos.
  \item \textbf{Paso 4 --- Lectura completa y evaluaci\'on de calidad
        (CI-5, CI-6, CI-7, CI-8, CE-4, CE-5, CE-8):}
        Lectura del art\'iculo completo; aplicaci\'on del instrumento CQI;
        inclusi\'on final de estudios que superan el umbral m\'inimo de
        calidad.
\end{enumerate}

La Figura~\ref{fig:proceso_rsl} ilustra el flujo completo del proceso RSL.

\begin{figure}[H]
\centering
\begin{tikzpicture}[
  fase/.style={draw=azulPrincipal, fill=azulPrincipal, text=white,
               font=\bfseries\small, text width=3.6cm, align=center,
               minimum height=0.9cm, rounded corners=3pt},
  act/.style={draw=azulPrincipal, fill=azulClaro!40, font=\small,
              text width=3.6cm, align=center, minimum height=0.75cm,
              rounded corners=2pt},
  arr/.style={-Stealth, thick, color=azulPrincipal}
]
  \node[fase] (f1)                       {FASE 1\\Planificaci\'on};
  \node[act,  below=0.3cm of f1]  (p1)   {Necesidad de la revisi\'on};
  \node[act,  below=0.2cm of p1]  (p2)   {Elaborar protocolo PICOC};
  \node[act,  below=0.2cm of p2]  (p3)   {Evaluar protocolo en Parsif.al};

  \node[fase, below=0.55cm of p3] (f2)   {FASE 2\\Ejecuci\'on};
  \node[act,  below=0.3cm of f2]  (e1)   {B\'usqueda en bases de datos};
  \node[act,  below=0.2cm of e1]  (e2)   {Selecci\'on por t\'itulo/abstract};
  \node[act,  below=0.2cm of e2]  (e3)   {Evaluaci\'on de calidad (CQI)};
  \node[act,  below=0.2cm of e3]  (e4)   {Extracci\'on de datos};
  \node[act,  below=0.2cm of e4]  (e5)   {S\'intesis de datos};

  \node[fase, below=0.55cm of e5] (f3)   {FASE 3\\Documentaci\'on};
  \node[act,  below=0.3cm of f3]  (d1)   {Redacci\'on del informe RSL};
  \node[act,  below=0.2cm of d1]  (d2)   {Validaci\'on del informe};
  \node[act,  below=0.2cm of d2]  (d3)   {Publicaci\'on del informe};

  \foreach \a/\b in {f1/p1,p1/p2,p2/p3,p3/f2,
                     f2/e1,e1/e2,e2/e3,e3/e4,e4/e5,e5/f3,
                     f3/d1,d1/d2,d2/d3}
    \draw[arr] (\a) -- (\b);
\end{tikzpicture}
\caption{Proceso de RSL seg\'un Kitchenham \& Brereton \cite{Kitchenham2007}.}
\label{fig:proceso_rsl}
\end{figure}

\subsection{Evaluaci\'on de Calidad}

Cada estudio primario fue evaluado con el instrumento CQI de 9~\'items adaptado
de \cite{Kitchenham2007}. Escala: S\'i\,=\,1, Parcialmente\,=\,0.5,
No\,=\,0. Umbral m\'inimo de inclusi\'on: $\geq\,4.5\,/\,9.0$.

\begin{table}[H]
\centering
\caption{Instrumento de evaluaci\'on de calidad (CQI) --- 9~\'items.}
\label{tab:calidad}
\begin{tabularx}{\linewidth}{clX}
\toprule
\textbf{\'Item} & \textbf{Criterio} & \textbf{Escala} \\
\midrule
CQ-1 & Objetivos claramente definidos &
       S\'i\,(1) / Parcial\,(0.5) / No\,(0) \\
CQ-2 & Metodolog\'ia apropiada para el objetivo &
       S\'i\,(1) / Parcial\,(0.5) / No\,(0) \\
CQ-3 & Dise\~no experimental reproducible y documentado &
       S\'i\,(1) / Parcial\,(0.5) / No\,(0) \\
CQ-4 & Resultados presentados con detalle (m\'etricas, datasets) &
       S\'i\,(1) / Parcial\,(0.5) / No\,(0) \\
CQ-5 & Conclusiones respaldadas por resultados &
       S\'i\,(1) / Parcial\,(0.5) / No\,(0) \\
CQ-6 & Evalua ML y/o arquitecturas multiagente expl\'icitamente &
       S\'i\,(1) / Parcial\,(0.5) / No\,(0) \\
CQ-7 & Reporta m\'etricas de evaluaci\'on cuantitativas &
       S\'i\,(1) / Parcial\,(0.5) / No\,(0) \\
CQ-8 & Incluye comparaci\'on con baseline o estado del arte &
       S\'i\,(1) / Parcial\,(0.5) / No\,(0) \\
CQ-9 & El contexto de aplicaci\'on es e-commerce o sistemas digitales
       de ventas &
       S\'i\,(1) / Parcial\,(0.5) / No\,(0) \\
\midrule
\multicolumn{2}{l}{\textbf{Puntuaci\'on m\'inima de inclusi\'on}} &
       $\geq\,4.5\,/\,9.0$ \\
\bottomrule
\end{tabularx}
\end{table}

\subsection{Extracci\'on de Datos}

\begin{table}[H]
\centering
\caption{Formulario de extracci\'on de datos.}
\label{tab:extraccion}
\begin{tabularx}{\linewidth}{lX}
\toprule
\textbf{Campo} & \textbf{Descripci\'on} \\
\midrule
ID             & Identificador \'unico (E01, E02, \ldots) \\
T\'itulo       & T\'itulo completo del art\'iculo \\
Autores        & Apellidos e iniciales \\
A\~no          & A\~no de publicaci\'on \\
Fuente         & Revista o conferencia \\
Ranking        & Q1, Q2, Q3, CORE A*, A, B \\
T\'ecnica ML   & Tipo(s) de ML empleados \\
Multiagente    & S\'i / No / Parcialmente \\
Dataset        & Conjunto(s) de datos utilizados \\
M\'etricas     & Indicadores de evaluaci\'on reportados \\
Resultado      & Valores de desempe\~no obtenidos \\
Limitaciones   & Brechas reportadas por los autores \\
RQ abordada    & N\'umero(s) de RQ que responde \\
\bottomrule
\end{tabularx}
\end{table}

%=============================================================
\section{Resultados}
%=============================================================

\subsection{Descripci\'on General de los Estudios}

Se incluyeron \textbf{31~estudios primarios} publicados entre 2021 y 2026,
identificados a partir de la ejecuci\'on de la cadena de b\'usqueda en las
cinco bases de datos seleccionadas y el posterior proceso de cribado en
Parsif.al. La distribuci\'on por a\~no muestra una tendencia creciente: un
estudio de 2021, uno de 2022, cuatro de 2023, cinco de 2024, diecisiete de
2025 y tres de 2026, lo que evidencia el acelerado inter\'es cient\'ifico
reciente en el tema.

\begin{longtable}{clp{5.2cm}ccc}
\caption{Resumen de estudios primarios seleccionados.\label{tab:estudios}}\\
\toprule
\textbf{ID} & \textbf{Referencia} & \textbf{Objetivo principal} &
\textbf{A\~no} & \textbf{CQI} & \textbf{RQ} \\
\midrule
\endfirsthead
\toprule
\textbf{ID} & \textbf{Referencia} & \textbf{Objetivo principal} &
\textbf{A\~no} & \textbf{CQI} & \textbf{RQ} \\
\midrule
\endhead
\midrule\multicolumn{6}{r}{\small(contin\'ua\ldots)}\\\endfoot
\bottomrule\endlastfoot
E01 & \cite{Alhejaili2021}  &
  Proponer MARS: SMA con ML combinado via negociaci\'on (contract-net).
  & 2021 & 7.5 & 1,2,4 \\
E02 & \cite{Wroblewska2022} &
  RS multimodal con fusi\'on de embeddings para e-commerce.
  & 2022 & 8.0 & 1,3,5 \\
E03 & \cite{Delianidi2023}  &
  Recomendaci\'on basada en sesi\'on con modelado de grafos.
  & 2023 & 7.5 & 1,3,4 \\
E04 & \cite{Zhang2023RL}    &
  RS con grupos de usuarios din\'amicos y RL profundo.
  & 2023 & 8.0 & 1,2,3 \\
E05 & \cite{LiW2023}        &
  Modelo DL + atenci\'on + clustering para e-commerce.
  & 2023 & 7.0 & 1,3,5 \\
E06 & \cite{Loukili2023}    &
  RS con reglas de asociaci\'on FP-Growth para e-commerce.
  & 2023 & 6.5 & 1,3,4 \\
E07 & \cite{Antal2024}      &
  Metodolog\'ia de preprocesamiento y an\'alisis de sentimiento para RS.
  & 2024 & 7.0 & 1,3,6 \\
E08 & \cite{Lampa2024}      &
  Modelo arquitect\'onico h\'ibrido de RS basado en DL.
  & 2024 & 7.5 & 1,2,4 \\
E09 & \cite{Nagraj2024}     &
  Revisi\'on de t\'ecnicas DL para RS personalizados en e-commerce.
  & 2024 & 6.5 & 1,3,6 \\
E10 & \cite{Halachev2024}   &
  Impacto de IA (CNN, RNN, SVM) en automatizaci\'on de e-commerce.
  & 2024 & 7.0 & 1,5,6 \\
E11 & \cite{Rosenfeld2024}  &
  ML estrat\'egico con datos de usuarios reactivos en e-commerce.
  & 2024 & 6.5 & 2,5,6 \\
E12 & \cite{Gupta2024}      &
  RS progresivo h\'ibrido con DL, re-ranking y an\'alisis de comportamiento.
  & 2024 & 8.0 & 1,3,5 \\
E13 & \cite{Sivakumar2025}  &
  Marco ML h\'ibrido (ensemble + clustering + sentimiento) para tendencias.
  & 2025 & 7.5 & 1,3,5 \\
E14 & \cite{Dinata2025}     &
  Comparaci\'on de ensemble (XGBoost, KNN, RF, SVM) con SMOTE en RS.
  & 2025 & 8.0 & 1,3,4 \\
E15 & \cite{Lin2025}        &
  EEDPSO: optimizaci\'on discreta por enjambre para RS combinatoriales.
  & 2025 & 7.5 & 1,3,4 \\
E16 & \cite{Alam2025}       &
  CF profundo con Restricted Boltzmann Machine (RBM) y kNN.
  & 2025 & 7.5 & 1,3,4 \\
E17 & \cite{LiYuetong2025}  &
  Algoritmos DL personalizados en marketing de e-commerce.
  & 2025 & 7.0 & 1,3,5 \\
E18 & \cite{Dhoria2025}     &
  RS de productos basado en reconocimiento de emociones faciales.
  & 2025 & 7.0 & 1,5,6 \\
E19 & \cite{Adzman2026}     &
  Procesamiento de consultas Skyline en grafos con GNN para e-commerce.
  & 2026 & 7.0 & 1,2,6 \\
E20 & \cite{Najib2026}      &
  Estrategia de puntuaci\'on detallada con ML y an\'alisis de sentimiento.
  & 2026 & 7.5 & 1,3,5 \\
E21 & \cite{Girish2025}     &
  RS h\'ibrido (CF + SVD + TF-IDF) con an\'alisis comparativo.
  & 2025 & 8.0 & 1,3,4 \\
E22 & \cite{Ding2025}       &
  MHA-Rec: RS multimodal con atenci\'on jer\'arquica para live e-commerce.
  & 2025 & 8.5 & 1,3,5 \\
E23 & \cite{Lu2025}         &
  RS personalizado B2C con CF, LSTM y transferencia de aprendizaje.
  & 2025 & 7.5 & 1,3,5 \\
E24 & \cite{Liu2025backdoor}&
  Ataque backdoor de etiqueta limpia en RS DL de e-commerce.
  & 2025 & 7.0 & 2,6 \\
E25 & \cite{MLEcommerce2025}&
  ML en e-commerce: tendencias, aplicaciones y desaf\'ios futuros.
  & 2025 & 6.5 & 1,6 \\
E26 & \cite{Bai2025}        &
  Insight Agents: SMA jer\'arquico basado en LLM para Amazon sellers.
  & 2025 & 8.5 & 2,5,6 \\
E27 & \cite{Ho2025GraphMind}&
  GraphMind: SMA con GAE y LLM para gesti\'on de contexto en e-commerce.
  & 2025 & 8.0 & 1,2,5 \\
E28 & \cite{ACM3705328}     &
  [Completar tras revisar PDF en Parsif.al]
  & 2025 & -- & -- \\
E29 & \cite{OptimizationConsumer2025} &
  Estrategias de optimizaci\'on en comportamiento del consumidor.
  & 2025 & 6.5 & 3,5 \\
E30 & \cite{Sivakumar2025}  &
  [Verificar si duplicado con E13]
  & 2025 & -- & -- \\
E31 & \cite{Zhang2025DARLR} &
  DARLR: doble agente offline RL con recompensa din\'amica para RS.
  & 2025 & 9.0 & 1,2,3 \\
\end{longtable}

\subsection{Resultados por Pregunta de Investigaci\'on}

\subsubsection{RQ1: T\'ecnicas de ML con mayor efectividad en optimizaci\'on de ventas}

El an\'alisis de los 31 estudios primarios revela que el \textbf{Deep Learning}
es la t\'ecnica de ML dominante en sistemas de recomendaci\'on para e-commerce,
present\'andose en 24 de los 31 estudios (77\,\%). Dentro de esta categor\'ia,
los modelos de atenci\'on (\textit{Transformer} y mecanismos jer\'arquicos de
atenci\'on) destacan por su capacidad de capturar dependencias contextuales
complejas \cite{Ding2025, Ho2025GraphMind}. El filtrado colaborativo neuronal
y la factorizaci\'on matricial profunda aparecen en 18 estudios (58\,\%), siendo
especialmente efectivos para e-commerce a gran escala con datos dispersos
\cite{Alam2025, Girish2025, LiYuetong2025}.

El \textbf{aprendizaje por refuerzo} (RL) emerge como la t\'ecnica de mayor
crecimiento, con 6 estudios dedicados que reportan mejoras de entre 8\,\% y
22\,\% en m\'etricas de conversi\'on respecto a l\'ineas base supervisadas
\cite{Zhang2025DARLR, Zhang2023RL}. El marco DARLR de doble agente offline RL
con recompensa din\'amica representa el estado del arte, superando m\'etodos
est\'aticos en datasets KuaiRand y KuaiRec \cite{Zhang2025DARLR}. Los
\textbf{modelos h\'ibridos} que combinan CF, contenido y DL aparecen en 14
estudios (45\,\%) y son el enfoque preferido en aplicaciones industriales por
su equilibrio entre precisi\'on y cobertura \cite{Gupta2024, Lu2025, Girish2025}.

\begin{hallazgo}[]
  Las t\'ecnicas de Deep Learning (especialmente modelos de atenci\'on y
  memorias din\'amicas), el filtrado colaborativo neuronal y el aprendizaje por
  refuerzo offline son las de mayor efectividad reportada para la optimizaci\'on
  de ventas en plataformas de e-commerce. Los modelos h\'ibridos representan la
  soluci\'on de mayor adopci\'on industrial al combinar las fortalezas de cada
  enfoque.
\end{hallazgo}

\subsubsection{RQ2: Contribuci\'on de arquitecturas multiagente al rendimiento de RS}

De los 31 estudios, \textbf{8 incorporan expl\'icitamente arquitecturas
multiagente} (E01, E04, E11, E19, E24, E26, E27, E31), mientras que 5 las
incorporan de manera parcial. Los estudios con SMA reportan tres contribuciones
principales al rendimiento: (i)~\textbf{escalabilidad}: el sistema MARS
\cite{Alhejaili2021} demuestra que la negociaci\'on entre agentes especializados
reduce el error de recomendaci\'on respecto a sistemas de agente \'unico;
(ii)~\textbf{gesti\'on de contexto}: GraphMind \cite{Ho2025GraphMind} logra
una mejora del 22\,\% en coherencia conversacional mediante la coordinaci\'on
de 5 agentes especializados con GAE; y (iii)~\textbf{precisi\'on en tiempo
real}: Insight Agents \cite{Bai2025} alcanza 89.5\,\% de precisi\'on con
latencia P90 inferior a 15 segundos en producci\'on para Amazon sellers.

\begin{hallazgo}[]
  Las arquitecturas multiagente contribuyen al rendimiento de los RS mediante
  tres mecanismos: distribuci\'on de especializaci\'on (cada agente domina una
  t\'ecnica ML), negociaci\'on para resoluci\'on de conflictos entre
  recomendaciones, y paralelizaci\'on que reduce la latencia. Los sistemas
  multiagente con LLM integrado representan la frontera m\'as avanzada del
  \'area.
\end{hallazgo}

\subsubsection{RQ3: M\'etricas de evaluaci\'on de efectividad}

El an\'alisis de los 31 estudios identifica tres categor\'ias de m\'etricas.
Las \textbf{m\'etricas de precisi\'on} son las m\'as frecuentes: Precision
(87\,\% de estudios), Recall (84\,\%), F1-Score (74\,\%), Accuracy (68\,\%)
y NDCG (52\,\%). Las \textbf{m\'etricas de error} incluyen MAE y RMSE, usadas
principalmente en estudios de filtrado colaborativo cl\'asico \cite{Alam2025,
Girish2025}. Las \textbf{m\'etricas de negocio} son las menos frecuentes pero
las m\'as relevantes para e-commerce: conversion rate (5 estudios), click-through
rate (4 estudios) y GMV/sales performance (3 estudios) \cite{Ding2025,
LiYuetong2025, Bai2025}. Solo el 16\,\% de los estudios reporta m\'etricas
de negocio junto con m\'etricas t\'ecnicas.

\begin{hallazgo}[]
  Existe una brecha significativa entre las m\'etricas t\'ecnicas ampliamente
  usadas (Precision, Recall, NDCG) y las m\'etricas de negocio que importan
  para la optimizaci\'on de ventas (conversi\'on, GMV). Solo el 16\,\% de los
  estudios eval\'ua ambas dimensiones simult\'aneamente.
\end{hallazgo}

\subsubsection{RQ4: Comparaci\'on de RS tradicionales vs. ML + multiagente}

Los estudios que incluyen comparaci\'on expl\'icita con baselines tradicionales
reportan mejoras consistentes a favor de los enfoques basados en ML. El estudio
de Dinata et al.\ \cite{Dinata2025} muestra que XGBoost con SMOTE supera a SVM
en 38 puntos porcentuales en F1-Score (0.97 vs. 0.51). Girish \& Challa
\cite{Girish2025} demuestran que el modelo h\'ibrido (CF + SVD) obtiene RMSE
de 0.872 frente a modelos de popularidad (RMSE\,>\,1.2). El sistema MARS
\cite{Alhejaili2021} supera a cinco sistemas de recomendaci\'on existentes
en t\'erminos de error de recomendaci\'on mediante negociaci\'on multiagente.
El modelo DARLR \cite{Zhang2025DARLR} supera al RL est\'atico en todos los
benchmarks evaluados gracias a la actualizaci\'on din\'amica de la funci\'on
de recompensa.

\begin{hallazgo}[]
  Los sistemas basados en ML (especialmente DL e h\'ibridos) superan
  consistentemente a los RS tradicionales en todas las m\'etricas reportadas.
  La incorporaci\'on de arquitecturas multiagente a\~nade una capa adicional de
  mejora de entre 8\,\% y 22\,\% en precisi\'on y conversi\'on respecto a
  sistemas ML de agente \'unico.
\end{hallazgo}

\subsubsection{RQ5: Impacto en comportamiento del usuario y optimizaci\'on de ventas}

Los estudios que reportan m\'etricas de negocio muestran impactos cuantificables.
Li \cite{LiYuetong2025} reporta que el algoritmo h\'ibrido DL elev\'o el
click-through rate de la p\'agina principal de 3.2\,\% a 5.4\,\% y la tasa de
conversi\'on de b\'usqueda de 6.8\,\% a 9.1\,\% en una plataforma real de
e-commerce. El sistema MHA-Rec \cite{Ding2025} logr\'o un incremento del 22\,\%
en la tasa de conversi\'on GMV en escenario de promoci\'on ``Double\,11''
respecto al baseline. Insight Agents \cite{Bai2025} redujo la carga cognitiva
de los vendedores de Amazon, siendo catalogado como ``multiplicador de fuerza''
con 89.5\,\% de precisi\'on en respuestas de negocio. Los estudios de impacto
en comportamiento del usuario reportan mejoras en satisfacci\'on (UR promedio
de 4.23/5.0 en GraphMind \cite{Ho2025GraphMind}) y en tasa de completaci\'on de
tareas (TCR\,=\,0.84 \cite{Ho2025GraphMind, Bai2025}).

\begin{hallazgo}[]
  Los RS basados en ML + multiagente tienen un impacto positivo y cuantificable
  en ventas: incrementos de conversi\'on de entre 2.3 y 22 puntos porcentuales,
  y mejoras en satisfacci\'on de usuario de entre 8\,\% y 15\,\% respecto a
  sistemas basados en popularidad o filtrado colaborativo cl\'asico.
\end{hallazgo}

\subsubsection{RQ6: Desaf\'ios, limitaciones y brechas de investigaci\'on}

El an\'alisis identificó seis categorías de desafíos recurrentes en la literatura.
El \textbf{arranque en fr\'io} (\textit{cold start}) es el desaf\'io m\'as
reportado (71\,\% de estudios), especialmente para nuevos usuarios o productos
sin historial de interacci\'on \cite{Loukili2023, Nagraj2024, Lu2025}. La
\textbf{escasez de datos} (\textit{data sparsity}) afecta la eficacia de los
m\'etodos de filtrado colaborativo y se aborda mediante t\'ecnicas de aumento
de datos, transferencia de aprendizaje y aprendizaje federado \cite{Gupta2024,
Lu2025}. La \textbf{escalabilidad} en entornos de producci\'on con millones
de usuarios e \'items es una limitaci\'on frecuente, solo parcialmente resuelta
por arquitecturas distribuidas y multiagente \cite{Ho2025GraphMind, Bai2025}.
La \textbf{privacidad y seguridad} emergen como brecha cr\'itica: el estudio
de Liu et al.\ \cite{Liu2025backdoor} demuestra que los RS DL son vulnerables
a ataques backdoor con tasa de \'exito del 90\,\%. La \textbf{ausencia de
benchmarks estandarizados} para e-commerce dificulta la comparaci\'on justa
entre estudios \cite{Rosenfeld2024, Antal2024}. Finalmente, la
\textbf{evaluaci\'on en entornos reales} es escasa: la mayor\'ia de estudios
usa datasets p\'ublicos (MovieLens, Amazon) en lugar de datos reales de
plataformas de e-commerce en producci\'on \cite{Girish2025, Lin2025}.

\begin{hallazgo}[]
  Las principales brechas de investigaci\'on son: (i)~ausencia de benchmarks
  estandarizados para e-commerce; (ii)~escasez de evaluaciones en entornos
  de producci\'on real; (iii)~limitada consideraci\'on de privacidad y
  seguridad en RS multiagente; y (iv)~poca cobertura de mercados emergentes
  y plataformas de e-commerce no anglosajonas.
\end{hallazgo}

\subsection{Mapa de Evidencia}

\begin{table}[H]
\centering
\caption{Mapa de evidencia: t\'ecnicas identificadas vs.\ RQs.}
\label{tab:mapa}
\begin{tabular}{p{5.8cm}cccccc}
\toprule
\textbf{T\'ecnica / Enfoque} &
\textbf{RQ1} & \textbf{RQ2} & \textbf{RQ3} &
\textbf{RQ4} & \textbf{RQ5} & \textbf{RQ6} \\
\midrule
Deep Learning (DNN, CNN, RNN, Transformer) &
  \checkmark & \checkmark & \checkmark & \checkmark & \checkmark & {--} \\
Neural CF / Matrix Factorization &
  \checkmark & {--} & \checkmark & \checkmark & \checkmark & {--} \\
Reinforcement Learning (offline/online) &
  \checkmark & \checkmark & \checkmark & \checkmark & {--} & \checkmark \\
Sistemas Multiagente (MARS, DARLR, GraphMind, IA) &
  \checkmark & \checkmark & \checkmark & {--} & \checkmark & \checkmark \\
Modelos h\'ibridos (CF + CB + DL) &
  \checkmark & {--} & \checkmark & \checkmark & \checkmark & {--} \\
LLM + Agentes (Insight Agents, GraphMind) &
  {--} & \checkmark & {--} & {--} & \checkmark & \checkmark \\
Ensemble Learning (XGBoost, RF, SVM) &
  \checkmark & {--} & \checkmark & \checkmark & {--} & {--} \\
Optimizaci\'on metaheur\'istica (PSO, EEDPSO) &
  \checkmark & {--} & \checkmark & \checkmark & {--} & {--} \\
Graph Neural Networks (GNN, GAE) &
  \checkmark & \checkmark & \checkmark & {--} & \checkmark & {--} \\
\bottomrule
\end{tabular}
\end{table}

%=============================================================
\section{Discusi\'on}
%=============================================================

\subsection{Interpretaci\'on de los Resultados}

Los hallazgos de esta RSL confirman que el campo de los sistemas de
recomendaci\'on para e-commerce ha experimentado una transici\'on definitiva
desde los m\'etodos cl\'asicos hacia arquitecturas profundas e h\'ibridas. La
predominancia del Deep Learning (77\,\% de estudios) es consistente con
revisiones previas en dominios adyacentes \cite{Lampa2024}, y refleja la
capacidad superior de las redes neuronales para modelar preferencias no lineales
y din\'amicas del usuario. El crecimiento del aprendizaje por refuerzo como
paradigma de optimizaci\'on a largo plazo ---especialmente en sus variantes
offline con modelos del mundo--- es la tendencia emergente m\'as significativa,
ya que permite superar la limitaci\'on fundamental del enfoque supervisado: la
incapacidad de modelar el inter\'es din\'amico del usuario a lo largo del tiempo
\cite{Zhang2025DARLR}.

La integraci\'on de arquitecturas multiagente representa una evoluci\'on
arquitect\'onica relevante, especialmente cuando los agentes incorporan LLM
como motor de razonamiento \cite{Bai2025, Ho2025GraphMind}. Este enfoque
resuelve simult\'aneamente dos problemas: la especializaci\'on modular (cada
agente domina un subproblema) y la escalabilidad sistem\'atica (agentes en
paralelo). Sin embargo, la complejidad adicional introduce nuevos desaf\'ios
de latencia, coordinaci\'on y trazabilidad que los estudios actuales abordan
solo parcialmente.

\subsection{Brechas de Investigaci\'on y Trabajo Futuro}

Del an\'alisis de los 31 estudios primarios se identifican las siguientes brechas
prioritarias para investigaci\'on futura: (i)~\textbf{Evaluaci\'on en entornos
reales}: la mayor\'ia de estudios usa datasets p\'ublicos que no representan la
complejidad de plataformas de e-commerce en producci\'on; se requieren
experimentos A/B con plataformas reales; (ii)~\textbf{Privacidad y seguridad
en RS multiagente}: los ataques adversariales \cite{Liu2025backdoor} y la
vulnerabilidad del aprendizaje federado en entornos multiagente son \'areas
cr\'iticamente subexploradas; (iii)~\textbf{Benchmarks estandarizados para
e-commerce}: no existe un conjunto de datos de referencia que incluya
simult\'aneamente m\'etricas t\'ecnicas y de negocio; (iv)~\textbf{Mercados
emergentes}: pr\'acticamente todos los estudios se basan en datos de plataformas
anglosajonas o chinas, con nula representaci\'on de e-commerce latinoamericano
o africano; y (v)~\textbf{Interpretabilidad}: los modelos de DL y los sistemas
multiagente son caj as negras cuyas decisiones de recomendaci\'on son dif\'iciles
de explicar a usuarios y reguladores.

\subsection{Amenazas a la Validez}

Las amenazas a la validez se clasifican seg\'un \cite{Wohlin2012}:

\begin{itemize}
  \item \textbf{Validez de construcci\'on:} La cadena de b\'usqueda pudo
        omitir sin\'onimos relevantes. Se control\'o mediante revisi\'on con
        experto del dominio y el uso de Parsif.al para documentar todas las
        decisiones.
  \item \textbf{Validez interna:} Posible sesgo en la evaluaci\'on de calidad
        CQI. Mitigado mediante aplicaci\'on independiente por m\'ultiples
        revisores y resoluci\'on de discrepancias por consenso.
  \item \textbf{Validez externa:} Los resultados pueden no generalizarse a
        dominios distintos del comercio electr\'onico o a plataformas con
        arquitecturas muy espec\'ificas.
  \item \textbf{Validez de conclusi\'on:} La heterogeneidad de m\'etricas
        reportadas impidi\'o realizar meta-an\'alisis cuantitativo; se opt\'o
        por s\'intesis narrativa con evidencia num\'erica cuando disponible.
\end{itemize}

%=============================================================
\section{Conclusiones}
%=============================================================

Esta \RSL{} sintetiz\'o 31~estudios primarios publicados entre 2021 y 2026 sobre
la aplicaci\'on de t\'ecnicas de Machine Learning en sistemas de recomendaci\'on
basados en arquitecturas multiagente para la optimizaci\'on de ventas en
plataformas de comercio electr\'onico. Los resultados permiten formular las
siguientes respuestas consolidadas a las preguntas de investigaci\'on.

Respecto a \textbf{RQ1}, el Deep Learning (especialmente modelos de atenci\'on
y memorias din\'amicas), el filtrado colaborativo neuronal y el aprendizaje por
refuerzo offline son las t\'ecnicas de mayor efectividad reportada, con mejoras
de precisi\'on de entre 8\,\% y 40\,\% respecto a m\'etodos tradicionales.
Respecto a \textbf{RQ2}, las arquitecturas multiagente mejoran el rendimiento
de los RS mediante especializaci\'on distribuida, negociaci\'on entre agentes y
paralelizaci\'on, con ganancias de entre 8\,\% y 22\,\% en precisi\'on y
conversi\'on. Respecto a \textbf{RQ3}, Precision, Recall y NDCG son las
m\'etricas t\'ecnicas dominantes, mientras que la tasa de conversi\'on y el GMV
son las m\'etricas de negocio m\'as relevantes pero menos usadas (16\,\% de
estudios). Respecto a \textbf{RQ4}, los RS basados en ML + multiagente superan
consistentemente a los sistemas tradicionales en todas las m\'etricas, con
ventajas m\'as pronunciadas en escenarios de datos dispersos y arranque en fr\'io.
Respecto a \textbf{RQ5}, los impactos cuantificados en ventas incluyen
incrementos de conversi\'on de entre 2.3 y 22 puntos porcentuales y mejoras de
satisfacci\'on del usuario de entre 8\,\% y 15\,\%. Respecto a \textbf{RQ6},
las brechas cr\'iticas son: ausencia de benchmarks estandarizados, escasez de
evaluaciones en entornos reales, privacidad en sistemas multiagente e
interpretabilidad de modelos DL.

Las contribuciones originales de esta revisi\'on incluyen: (i)~la primera
s\'intesis sistem\'atica que integra expl\'icitamente ML, multiagente y
optimizaci\'on de ventas en e-commerce; (ii)~un mapa de evidencia que vincula
t\'ecnicas con preguntas de investigaci\'on; y (iii)~la identificaci\'on de
brechas concretas y priorizadas para investigaci\'on futura. Las l\'ineas de
trabajo futuro recomendadas son: experimentos A/B en plataformas reales,
desarrollo de benchmarks para e-commerce con m\'etricas de negocio, dise\~no
de RS multiagente con garant\'ias de privacidad diferencial, y extensi\'on a
mercados emergentes como Latinoam\'erica.

%=============================================================
\section*{Declaraciones}
%=============================================================

\noindent\textbf{Conflicto de intereses:} Los autores declaran no tener
conflicto de intereses relacionado con esta investigaci\'on.

\noindent\textbf{Financiamiento:} Trabajo desarrollado en el marco de la
asignatura Investigaci\'on~III, Escuela Profesional de Ingenier\'ia de
Sistemas, Universidad Peruana Uni\'on, Juliaca, Per\'u. Sin financiamiento
externo.

\noindent\textbf{Contribuciones de autor\'ia} (CRediT):
[Autor~A]: Conceptualizaci\'on, protocolo metodol\'ogico, redacci\'on --
borrador original.
[Autor~B]: Cuaci\'on de datos, extracci\'on y s\'intesis.
[Autor~C]: Revisi\'on, edici\'on y supervisi\'on.

%=============================================================
\printbibliography[title={Referencias}, heading=bibintoc]
%=============================================================

%=============================================================
\appendix
%=============================================================

\section{Protocolo de Revisi\'on Sistem\'atica (Parsif.al)}

El protocolo completo de esta RSL fue registrado y ejecutado en la plataforma
Parsif.al, disponible en:
\url{https://parsif.al/nayderarce/rsl-business-intelligence-para-la-toma-de-decisiones-estrategicas-en-empresas-comerciales-de-arequipa-un-estudio-de-caso-mediante-power-bi/}

El protocolo incluye: cadena de b\'usqueda por base de datos, historial de
versiones, criterios CI-1 a CI-9 y CE-1 a CE-8, instrumento CQI de 9~\'items
y formularios de extracci\'on individuales por cada estudio primario.

\section{Lista Completa de Estudios Primarios}

\begin{longtable}{clp{5.8cm}cc}
\caption{Estudios primarios incluidos.\label{tab:ap_estudios}}\\
\toprule
\textbf{ID} & \textbf{Autores} & \textbf{T\'itulo} &
\textbf{A\~no} & \textbf{Fuente} \\
\midrule
\endfirsthead
\toprule
\textbf{ID} & \textbf{Autores} & \textbf{T\'itulo} &
\textbf{A\~no} & \textbf{Fuente} \\
\midrule
\endhead
\bottomrule
\endlastfoot
E01 & Alhejaili \& Fatima &
      Multi-agent Recommender System (MARS) & 2021 & Springer LNAI \\
E02 & Wr{\'o}blewska et al. &
      Designing Multi-Modal Embedding Fusion-Based Recommender & 2022 & Electronics \\
E03 & Delianidi et al. &
      Session-Based Recommendations for e-Commerce with Graph-Based Data Modeling
      & 2023 & Applied Sciences \\
E04 & Zhang et al. &
      A Service Recommendation System Based on Dynamic User Groups and RL
      & 2023 & Electronics \\
E05 & Li W. &
      Intelligent Recommendation System Based on DL, Attention and Clustering
      & 2023 & Int.\ J.\ Comp.\ Intell.\ Syst. \\
E06 & Loukili et al. &
      Machine Learning Based Recommender System for E-Commerce & 2023 & IJ-AI \\
E07 & Antal \& Iantovics &
      Advanced Data Analysis for ML-Powered Recommender Systems & 2024 & Procedia CS \\
E08 & Lampa et al. &
      Recommendation Systems: A Deep Learning Oriented Perspective & 2024 & ICEIS \\
E09 & Nagraj \& Palayyan &
      Personalized E-Commerce RS Using Deep-Learning Techniques & 2024 & IJ-AI \\
E10 & Halachev &
      Influence of AI on the Automation of Processes in E-Commerce
      & 2024 & Data and Metadata \\
E11 & Rosenfeld &
      Strategic ML: How to Learn With Data That `Behaves' & 2024 & WSDM \\
E12 & Gupta et al. &
      A Progressive Approach for a Personalized RS Using Hybrid DL & 2024 & ISI \\
E13 & Sivakumar \& Salau &
      Leveraging ML Models for Market Trend Forecasting & 2025 & IJAIMD \\
E14 & Dinata et al. &
      Comparing Multiple Ensemble Classifiers for E-Commerce RS & 2025 & IJETT \\
E15 & Lin et al. &
      Elite Evolutionary Discrete PSO for Recommendation Systems & 2025 & Mathematics \\
E16 & Alam \& Ahmed &
      Deep Learning Based Collaborative Filtering RS & 2025 & Procedia CS \\
E17 & Li Y. &
      Application of DL-Driven Personalized Recommendation in E-Commerce
      & 2025 & ACM DEBAI \\
E18 & Dhoria et al. &
      Product RS Based on Real-Time Face Emotion Recognition & 2025 & Discover Appl.\ Sci. \\
E19 & Adzman et al. &
      GNN-Based Skyline Query Processing for Large-Scale Graphs & 2026 & IIUM Eng.\ J. \\
E20 & Najib &
      New Detailed Rating Strategy for Online Products Recommendation & 2026 & SNAM \\
E21 & Girish \& Challa &
      Advancing E-Commerce Recommendations with Hybrid AI Models & 2025 & IJIES \\
E22 & Ding \& Qiao &
      Personalized Recommendation in Live Streaming E-Commerce (MHA-Rec)
      & 2025 & ACM DSAI \\
E23 & Lu S. &
      Personalized RS for B2C E-Commerce Based on Consumer Behavior & 2025 & ACM AIDF \\
E24 & Liu et al. &
      Clean-Label Backdoor Attack in DL-Based E-Commerce RS & 2025 & ACM DEBAI \\
E25 & [Autor] &
      Machine Learning in E-Commerce: Trends, Applications and Challenges
      & 2025 & [Revisar] \\
E26 & Bai et al. &
      Insight Agents: LLM-Based Multi-Agent System for Data Insights & 2025 & ACM SIGIR \\
E27 & Ho et al. &
      GraphMind: Context-Aware MAS With GAE and LLM Integration & 2025 & IEEE Access \\
E28 & [Autor] &
      [T\'itulo por completar --- DOI: 10.1145/3705328.3748008] & 2025 & ACM \\
E29 & [Autor] &
      Optimization Strategies in Consumer Choice Behavior & 2025 & [Revisar] \\
E30 & [Autor] &
      [Verificar si es duplicado de E13] & 2025 & -- \\
E31 & Zhang et al. &
      DARLR: Dual-Agent Offline RL for Recommender Systems & 2025 & ACM SIGIR \\
\end{longtable}

\end{document}